\documentclass[12pt,a4paper]{article}
\usepackage[T2A]{fontenc}
\usepackage[utf8]{inputenc}
\usepackage[russian]{babel}
\usepackage{graphicx}
\usepackage[a4paper,margin=2.2cm]{geometry}
\usepackage{parskip}

\begin{document}

\begin{center}
    {\LARGE \textbf{Мария Склодовская-Кюри}}\\[0.5cm]
    \includegraphics[width=0.32\textwidth]{marie_curie_school_portrait_simple.png}
\end{center}

Мария Склодовская-Кюри - известная учёная польского происхождения.
Она жила и работала во Франции.
Мария изучала физику и химию.
Она стала знаменитой благодаря исследованиям радиоактивности.
Вместе с мужем Пьером Кюри она открыла полоний и радий.
За свои открытия она получила две Нобелевские премии.
Мария Кюри была первой женщиной, получившей Нобелевскую премию.
Сегодня её считают одной из самых важных женщин в истории науки.

\end{document}
